\documentclass[12pt]{article}
\usepackage[T1]{fontenc}
\usepackage[utf8]{inputenc}
\usepackage{lmodern}
\usepackage{textcomp}
\usepackage{enumitem}
\usepackage{geometry}
\usepackage{graphicx}
\usepackage{amsmath, amssymb}
\geometry{margin=1in}
\begin{document}
\begin{center}
Chemistry - Graduate Level Assessment 3 \
Total Marks: 40 \
Answer all questions.
\end{center}

\section*{Multiple Choice (10 marks)}
\begin{enumerate}[label=\arabic*.]
\item Calculate the volume of 0.3 N base necessary to neutralize 3 liters of 0.01 N nitric acid.
\begin{enumerate}[label=(\alph*)]
    \item 0.3 liters
    \item 2 liters
    \item 0.2 liters
    \item 3 liters
    \item 0.03 liters
\end{enumerate}
\item How many milliliters of 0.250 M KOH does it take to neutralize completely 50.0 mL of 0.150 M H$_{3}$PO$_{4}$?
\begin{enumerate}[label=(\alph*)]
    \item 75.0 mL
    \item 90.0 mL
    \item 60.0 mL
    \item 120 mL
    \item 30.0 mL
\end{enumerate}
\item Compute the missing mass of (^19 \_9)F (atomic weight = 18.9984) when the masses of proton, neutron and electron are 1.00728, 1.00867 and 0.000549 respectively.
\begin{enumerate}[label=(\alph*)]
    \item 0.1645 amu
    \item 0.1987 amu
    \item 0.1312 amu
    \item 0.1588 amu
    \item 0.1756 amu
\end{enumerate}
\item 4 liters of octane gasoline weigh 3.19 kg. Calculate what volume of air is required for its complete combustion at S.T.P.
\begin{enumerate}[label=(\alph*)]
    \item 22,102.28 liters air
    \item 58,650.00 liters air
    \item 49,350.25 liters air
    \item 65,292.15 liters air
    \item 39,175.40 liters air
\end{enumerate}
\item Without referring to a table, place the following hydro-carbons in order of increasing boiling points. (a) methane(d)neopentane (b) n-hexane(e) 2,3-dimethylbutane (c) n-undecane
\begin{enumerate}[label=(\alph*)]
    \item 2,3-dimethylbutane \textless{} neopentane \textless{} n-hexane \textless{} n-undecane \textless{} methane
    \item methane \textless{} 2,3-dimethylbutane \textless{} neopentane \textless{} n-hexane \textless{} n-undecane
    \item n-undecane \textless{} n-hexane \textless{} neopentane \textless{} 2,3-dimethylbutane \textless{} methane
    \item neopentane \textless{} 2,3-dimethylbutane \textless{} methane \textless{} n-undecane \textless{} n-hexane
    \item 2,3-dimethylbutane \textless{} methane \textless{} n-hexane \textless{} neopentane \textless{} n-undecane
\end{enumerate}
\end{enumerate}

\section*{True / False (10 marks)}
\textit{Answer True or False}
\begin{enumerate}[label=\arabic*.]
\setcounter{enumi}{5}
\item True or False: The entropy change for the vaporization of water is approximately 16.2 J mol-1 K-1.
\item True or False: The maximum wavelength of the Earth's blackbody radiation at a mean surface temperature of 288 K is approximately 0.7 x 10^\{-5\} m.
\item True or False: BF 3 possesses an octet of electrons around the central boron atom, fulfilling the octet requirement.
\item True or False: Oxygen diffuses faster than hydrogen because it is a lighter gas compared to hydrogen.
\item True or False: The molarity of the H$_{2}$SO$_{4}$ solution is calculated to be 0.265 M based on the neutralization reaction with NaOH.
\end{enumerate}

\section*{Long Form Response (20 marks)}
\textit{Answer each question with a clear and complete explanation.}
\begin{enumerate}[label=\arabic*.]
\setcounter{enumi}{10}
\item A chemist discovered a new compound and observed a freezing point depression of water to -1.22^\ensuremath{\circC} with 2.00 grams of the compound. Calculate its molecular weight and discuss the principles involved in this determination.
\item Calculate the number of molecules in 1.1 \ensuremath{\times} $10^8$ kg of phosgene (CCl 2 O) used in World War I, and discuss the implications of chemical warfare on human health and ethics.
\end{enumerate}

\vfill
\noindent\textit{End of Paper}
\end{document}